% 原试卷 PDF 第44～45页。
\setchoices{label=(\Alph*)}

\begin{question}[index=1]
$\displaystyle\lim_{n\to\infty}\frac{1}{n}\sum_{i=1}^{n}
\int_1^{\frac{i}{n}}\frac{\sin x}{x}\dd x=$
\begin{choices}[columns=2]
\item $1-2\cos 1$．
\item $\cos 1-1$．
\item $1-2\sin 1$．
\item $\sin 1-1$．
\end{choices}
\end{question}
\QuestionBreak

\begin{question}[index=2]
设函数 $f(x)$ 与 $g(x)$ 满足 $f(x_0)=g(x_0)$，$f'(x_0)<g'(x_0)$，则必存在 $\delta>0$，使得
\begin{choices}[columns=1]
\item 当 $x\in(x_0-\delta,x_0)\cup(x_0,x_0+\delta)$ 时，有 $f(x)<g(x)$．
\item 当 $x\in(x_0-\delta,x_0)\cup(x_0,x_0+\delta)$ 时，有 $f(x)>g(x)$．
\item 当 $x\in(x_0-\delta,x_0)$ 时，有 $f(x)>g(x)$，当 $x\in(x_0,x_0+\delta)$ 时，有 $f(x)<g(x)$．
\item 当 $x\in(x_0-\delta,x_0)$ 时，有 $f(x)<g(x)$，当 $x\in(x_0,x_0+\delta)$ 时，有 $f(x)>g(x)$．
\end{choices}
\end{question}
\QuestionBreak

\begin{question}[index=3]
设方程 $y''+ay'+y=(x+1)\mathrm e^{bx}$ 的每一个解 $y=y(x)$ 在区间 $[0,+\infty)$ 上均有界，则实数 $a,b$ 满足
\begin{choices}[columns=4]
\item $a\geq 0,\ b<0$．
\item $a\geq 0,\ b>0$．
\item $a>0,\ b>0$．
\item $a<0,\ b>0$．
\end{choices}
\end{question}
\QuestionBreak

\begin{question}[index=4]
若 $\displaystyle\frac{(y^2+2xy+ax^2)\dd x-(x^2+2xy+by^2)\dd y}{(x^2+y^2)^2}$ 为某函数 $u(x,y)$ 的全微分，则
\begin{choices}[columns=2]
\item $a=1,\ b=1$．
\item $a=1,\ b=-1$．
\item $a=-1,\ b=1$．
\item $a=-1,\ b=-1$．
\end{choices}
\end{question}
\QuestionBreak

\begin{question}[index=5]
设3阶矩阵 $\symbfit A,\symbfit B$ 满足 $r(\symbfit B\symbfit A)<r(\symbfit A\symbfit B)$，对于以下结论\par
① $\symbfit A\symbfit B\symbfit x=\symbf{0}$ 与 $\symbfit B\symbfit A\symbfit x=\symbf{0}$ 有非零公共解；\par
② $\symbfit A\symbfit B\symbfit A\symbfit x=\symbf{0}$ 与 $\symbfit B\symbfit A\symbfit B\symbfit x=\symbf{0}$ 有非零公共解．\par
正确的说法是
\begin{choices}[columns=2]
\item ①正确，②正确．
\item ①正确，②错误．
\item ①错误，②正确．
\item ①错误，②错误．
\end{choices}
\end{question}
\QuestionBreak

\begin{question}[index=6]
以下两个矩阵，可用同一可逆矩阵 $\symbfit P$ 相似对角化的是
\begin{choices}[columns=2,label-pos=left]
\item $\begin{pmatrix}1&1\\1&0\end{pmatrix},\begin{pmatrix}0&1\\1&1\end{pmatrix}$．
\item $\begin{pmatrix}1&1\\1&-1\end{pmatrix},\begin{pmatrix}-1&1\\1&1\end{pmatrix}$．
\item $\begin{pmatrix}0&1\\1&1\end{pmatrix},\begin{pmatrix}-1&1\\1&0\end{pmatrix}$．
\item $\begin{pmatrix}0&1\\1&-1\end{pmatrix},\begin{pmatrix}-1&1\\1&0\end{pmatrix}$．
\end{choices}
\end{question}
\QuestionBreak

\begin{question}[index=7]
已知 $\symbfit A=\begin{pmatrix}2&a-6&-2\\a+3&1&-4\\4&2&-4\end{pmatrix}$，则 $a=1$ 是存在两个不同的3阶矩阵 $\symbfit B$ 与 $\symbfit C$，使得 $\symbfit A\symbfit B=\symbfit A\symbfit C$ 成立的
\begin{choices}[columns=2]
\item 充分必要条件．
\item 充分非必要条件．
\item 必要非充分条件．
\item 既非充分也非必要条件．
\end{choices}
\end{question}
\QuestionBreak

\begin{question}[index=8]
设总体 $X$ 的概率密度为
$f(x)=\begin{cases}\dfrac{2x}{\alpha^2},&0\leq x\leq\alpha,\\0,&\text{其他，}\end{cases}$
其中 $\alpha(\alpha>1)$ 是未知参数，$X_1,X_2,\ldots,X_n$ 是来自总体 $X$ 的简单随机样本，记 $p=P\{0<X<\sqrt{\alpha}\}$，
$X_{(1)}=\min\{X_1,X_2,\ldots,X_n\}$，
$X_{(n)}=\max\{X_1,X_2,\ldots,X_n\}$，则 $p$ 的最大似然估计量 $\hat p$ 为
\begin{choices}[columns=4,label-pos=left]
\item $\sqrt{X_{(1)}}$．
\item $\sqrt{X_{(n)}}$．
\item $\dfrac{1}{X_{(1)}}$．
\item $\dfrac{1}{X_{(n)}}$．
\end{choices}
\end{question}
\QuestionBreak

\begin{question}[index=9]
已知随机变量 $X$ 的概率密度为
$f(x)=\begin{cases}
\dfrac{1}{C\sqrt{2\pi}\sigma}\mathrm e^{-\frac{1}{2\sigma^2}(x-\mu)^2},&x\geq 0,\\
0,&x<0,
\end{cases}$
其中 $C>0$，$\sigma>0$．记标准正态分布函数为 $\Phi(t)$，概率密度为 $\varphi(t)$，则 $C$ 的值与 $X$ 的期望 $EX$ 分别为
\begin{choices}[columns=2,label-pos=left]
\item $\Phi\left(\dfrac{\mu}{\sigma}\right),\
\mu+\varphi\left(\dfrac{\mu}{\sigma}\right)\dfrac{\sigma}{\Phi\left(\dfrac{\mu}{\sigma}\right)}$．
\item $\Phi\left(\dfrac{\mu}{\sigma}\right),\
\mu+\varphi\left(\dfrac{\mu}{\sigma}\right)\dfrac{\sigma\mu}{\Phi\left(\dfrac{\mu}{\sigma}\right)}$．
\item $\Phi\left(\dfrac{\mu}{\sigma}\right)+\mu,\
\varphi\left(\dfrac{\mu}{\sigma}\right)\dfrac{\sigma}{\Phi\left(\dfrac{\mu}{\sigma}\right)}$．
\item $\Phi\left(\dfrac{\mu}{\sigma}\right)+\mu,\
\varphi\left(\dfrac{\mu}{\sigma}\right)\dfrac{\sigma\mu}{\Phi\left(\dfrac{\mu}{\sigma}\right)}$．
\end{choices}
\end{question}
\QuestionBreak

\begin{question}[index=10]
设随机变量 $X,Y,Z$ 满足 $DX=DZ$，$DY=4DX$，且 $\rho_{XY}=-1$，$\rho_{XZ}=\dfrac12$，则 $X$ 与 $Y+Z$ 的相关系数为
\begin{choices}[columns=4,label-pos=left]
\item $\dfrac{\sqrt3}{2}$．
\item $-\dfrac{\sqrt3}{2}$．
\item $\dfrac34$．
\item $-\dfrac34$．
\end{choices}
\end{question}
